\documentclass[11pt,a4paper]{article}
\usepackage[margin=1in]{geometry}
\usepackage{times}
\usepackage{latexsym}
\usepackage{amsmath,amssymb}
\usepackage{booktabs}
\usepackage{graphicx}
\usepackage{multirow}
\usepackage{xcolor}
\usepackage{url}
\usepackage{hyperref}
\usepackage{microtype}
\usepackage{tabularx}
\usepackage{natbib}
\bibliographystyle{plainnat}

\newcommand{\ghat}{\hat{\gamma}}
\newcommand{\contamcal}{\textsc{ContamCal}}

\title{ContamCal: Calibrating LLM Benchmark Scores via\\Contamination-Augmented Item Response Theory}

\author{Anonymous}

\begin{document}
\maketitle

\begin{abstract}
Benchmark contamination---where training data overlaps with evaluation items---undermines the reliability of LLM evaluation, yet existing detection methods flag contamination without quantifying its effect on scores. We present \contamcal{}, a two-stage framework that bridges contamination detection and score calibration. Stage~1 computes item-level exposure scores via embedding similarity and fuses them with Reciprocal Rank Fusion. Stage~2 fits a contamination-augmented 4-parameter logistic IRT model with per-benchmark contamination coefficients~$\gamma_k$, producing calibrated scores with Bayesian credible intervals. Through controlled SFT injection experiments across Qwen2.5 (0.5B--7B) and Llama-3.2-1B at four dosage levels on four benchmarks, we establish three findings: (1)~embedding-based detection is the only reliable signal across all scales (AUC${}>{}$0.95), while Min-K++ and Self-Critique remain near chance---a method-specific rather than scale-dependent failure; (2)~contamination effects are benchmark-dependent, with ARC-Challenge exhibiting significant forgetting ($\ghat=-1.46$, 95\% CI $[-2.51,-0.41]$) while GSM8K shows no significant effect; (3)~Fisher Information analysis establishes identifiability thresholds (structural: AUC${\geq}$0.55; practical: AUC${\geq}$0.70) and confirms that coverage gaps at $N{=}16$ models are attributable to insufficient sample size (${>}70$\%), not model misspecification. Min-K++ exhibits directional reversal at high contamination (Cohen's $d=-1.45$ to $-1.98$), actively misleading rather than merely failing. We release controlled contamination models and evaluation code to support reproducible contamination research.
\end{abstract}

%============================================================
\section{Introduction}
\label{sec:intro}
%============================================================

Benchmark scores are the primary currency of LLM evaluation, yet they are increasingly debased by training data contamination. Rephrased test items can inflate MMLU accuracy by 8--15 points while evading standard $n$-gram decontamination \cite{xu2024benchmark}. The problem is worsening: \citet{balloccu2024leak} document contamination rates exceeding 90\% on popular benchmarks for recent models. Most critically, \citet{xu2025emperor} systematically evaluated 20 mitigation strategies and found \emph{none} reliably eliminates contamination effects---motivating a fundamental shift from prevention to post-hoc calibration.

Existing approaches address pieces of the puzzle but not the whole. Detection methods---Min-K\%++ \cite{shi2024detecting}, ConStat \cite{dekoninck2024constat}, Self-Critique---can flag likely contamination but produce no calibrated score answering ``what would this model score if it had never seen the test data?'' IRT-based evaluation methods \cite{lalor2024aaai,polo2024tinybenchmarks} model item difficulty psychometrically but ignore contamination. DCR \cite{oren2024dcr} provides adjusted accuracy via fuzzy matching but lacks uncertainty quantification and assumes a specific contamination mechanism.

\contamcal{} reframes benchmark contamination as a measurable covariate rather than a binary condition to eliminate. Grounded in educational measurement's item compromise IRT tradition \cite{eckerly2017detecting,sinharay2017detection} and explanatory IRT \cite{deboeck2004explanatory}, \contamcal{} augments a 4-parameter logistic IRT model with a contamination exposure covariate $\gamma_k$, estimated per benchmark~$k$. This allows contamination to manifest differently across benchmarks---inflation on knowledge-intensive tasks, \emph{forgetting} on reasoning tasks---while providing calibrated scores with Bayesian credible intervals that honestly reflect epistemic uncertainty.

Our contributions are:
\begin{enumerate}
\item \textbf{Signal effectiveness characterization.} A systematic mapping of detection signal $\times$ scale $\times$ dosage reveals that embedding similarity is the only reliable signal across 0.5B--3B scales (AUC$>$0.95), while Min-K++ and Self-Critique remain at chance level. Min-K++ exhibits directional reversal at high contamination (Cohen's $d = -1.45$ to $-1.98$), actively \emph{misleading} rather than merely failing---a cautionary finding for the community (\S\ref{sec:signals}).

\item \textbf{Contamination-augmented IRT calibration.} Per-benchmark $\gamma_k$ coefficients capture benchmark-dependent contamination effects: ARC-Challenge shows significant forgetting ($\ghat=-1.46$, CI excludes zero) while GSM8K shows no significant effect ($\ghat=0.19$, CI includes zero). This is a proof-of-concept that IRT$+\gamma$ captures directional effects, not a claim of MAE superiority (\S\ref{sec:irt}).

\item \textbf{Fisher Information identifiability analysis.} A 120-condition simulation sweep establishes structural (AUC$\geq$0.55) and practical (AUC$\geq$0.70) identifiability thresholds, and confirms that coverage gaps at $N{=}16$ are dominated by insufficient sample size ($>$70\%), providing actionable guidance for practitioners (\S\ref{sec:identifiability}).
\end{enumerate}

%============================================================
\section{Related Work}
\label{sec:related}
%============================================================

\paragraph{Contamination detection.}
Detection methods range from string matching and $n$-gram overlap \cite{dodge2021documenting} to membership inference approaches. Min-K\%++ \cite{shi2024detecting} uses token-level perplexity differences; ConStat \cite{dekoninck2024constat} compares original and perturbed item performance; Self-Critique elicits model self-assessment. These methods detect but do not calibrate. Our work treats detection outputs as \emph{inputs} to a calibration model.

\paragraph{IRT for LLM evaluation.}
Item Response Theory has been applied to LLM evaluation by \citet{lalor2024aaai}, TinyBenchmarks \cite{polo2024tinybenchmarks}, ATLAS, and Growing Pains. HiBayES provides hierarchical Bayesian evaluation. None model contamination as a covariate. Our $\gamma$-augmented IRT extends explanatory IRT \cite{deboeck2004explanatory} to the contamination setting, drawing on the item compromise IRT literature \cite{eckerly2017detecting,sinharay2017detection}.

\paragraph{Contamination mitigation.}
\citet{xu2025emperor} tested 20 mitigation strategies (rephrasing, option shuffling, etc.) and found none effective. DCR \cite{oren2024dcr} provides adjusted accuracy via fuzzy matching but lacks uncertainty quantification. LiveBench and similar dynamic benchmarks \cite{white2024livebench} avoid contamination via continuous curation but require ongoing effort. \contamcal{} complements these by providing post-hoc calibration with uncertainty for static benchmarks.

\paragraph{Membership inference and data attribution.}
Min-K++ reversal connects to broader MIA literature \cite{shokri2017membership,carlini2022extracting}: token-level memorization signals can invert under distributional shift from fine-tuning. Data attribution methods \cite{pruthi2020tracin,park2023trak} offer potential future signal sources.

%============================================================
\section{Method}
\label{sec:method}
%============================================================

\contamcal{} operates in two stages: (1)~item-level contamination exposure scoring via detection signal fusion, and (2)~contamination-augmented IRT calibration.

\subsection{Stage 1: Exposure Score Computation}
\label{sec:stage1}

For each test item $i$ and model $m$, we compute detection signals:
\begin{itemize}
\item \textbf{S1 (Embedding similarity).} Cosine similarity between the item embedding and nearest training data embedding, computed via a frozen sentence encoder. Training-agnostic: does not require model internals beyond embeddings.
\item \textbf{S2 (Min-K++).} Token-level perplexity ratio between the evaluated model and a reference model \cite{shi2024detecting}.
\item \textbf{S4 (Self-Critique).} Model self-assessment of item familiarity via prompted generation.
\end{itemize}

Signals are fused via Reciprocal Rank Fusion \cite{cormack2009rrf}:
\begin{equation}
\text{RRF}(i) = \sum_{s \in \mathcal{S}} \frac{1}{k + \text{rank}_s(i)}
\end{equation}
where $k=60$ is the standard smoothing constant. RRF is chosen over learned aggregation to avoid meta-circularity: a learned aggregator's training data may itself be contaminated. We show empirically (\S\ref{sec:rrf_results}) that RRF degrades gracefully to dominant-signal mode when constituent signals fail, rather than introducing noise.

\subsection{Stage 2: Contamination-Augmented 4PL IRT}
\label{sec:stage2}

We model the probability that model~$m$ answers item~$i$ on benchmark~$k$ correctly as:
\begin{equation}
P(Y_{mi} = 1) = c_i + (d_i - c_i) \cdot \sigma\bigl(a_i(\theta_m + \gamma_k \cdot e_m - b_i)\bigr)
\label{eq:4pl_gamma}
\end{equation}
where $\theta_m$ is model ability, $a_i$ is item discrimination, $b_i$ is item difficulty, $c_i$ and $d_i$ are lower and upper asymptotes, $e_m$ is model $m$'s exposure score (from Stage~1), and $\gamma_k$ is the per-benchmark contamination coefficient. $\sigma(\cdot)$ is the logistic function.

The key departure from standard IRT: $\gamma_k$ is estimated \emph{per benchmark}, allowing contamination to manifest as inflation ($\gamma_k > 0$), forgetting ($\gamma_k < 0$), or null effect ($\gamma_k \approx 0$). This is structurally equivalent to differential item functioning (DIF) between exposure groups \cite{holland1993dif,zumbo1999dif}.

\paragraph{Bayesian estimation.} We fit the model in Stan with weakly informative priors: $\theta_m \sim \mathcal{N}(0,1)$, $\gamma_k \sim \mathcal{N}(0,1)$, $a_i \sim \text{LogNormal}(0,0.5)$, $b_i \sim \mathcal{N}(0,2)$. Inference uses NUTS with 4 chains, 1000 warmup, 1000 sampling iterations. Convergence is assessed via $\hat{R} < 1.05$ and ESS $> 100$.

\paragraph{Calibrated score.} The decontaminated score for model $m$ on benchmark $k$ is:
\begin{equation}
\hat{s}_{mk} = \frac{1}{|\mathcal{I}_k|}\sum_{i \in \mathcal{I}_k} P(Y_{mi}=1 \mid \theta_m, \gamma_k=0, \hat{a}_i, \hat{b}_i, \hat{c}_i, \hat{d}_i)
\end{equation}
with credible intervals derived from posterior samples.

\subsection{Fisher Information and Identifiability}
\label{sec:fisher}

The Fisher Information for $\gamma_k$ in our model is:
\begin{equation}
I(\gamma_k) \propto N_{\text{models}} \cdot \text{Var}(e) \cdot f(\text{AUC})
\end{equation}
where $f(\text{AUC})$ captures the information content of exposure scores. When signal AUC$\to 0.5$ (chance), exposure scores become uninformative and $I(\gamma_k) \to 0$ regardless of $N$. Conversely, $N_{\text{models}}$ is a multiplicative factor that AUC cannot compensate---even perfect signals require sufficient models for reliable estimation.

%============================================================
\section{Experimental Setup}
\label{sec:setup}
%============================================================

\subsection{Controlled Contamination Injection}

We inject contamination via supervised fine-tuning (SFT) with LoRA on Qwen2.5 at three scales (0.5B, 1.5B, 3B) and four dosage levels: d000 (clean), d005 (5\% test items), d025 (25\%), d100 (100\%). Contamination data consists of semantically equivalent paraphrases of benchmark test items. We evaluate on four benchmarks: MMLU (14K items), ARC-Challenge (1.2K), GSM8K (1.3K), and HumanEval (164). For scale bridging, we additionally train Qwen2.5-7B at five dosage levels. For cross-family validation, we train Llama-3.2-1B at three dosage levels (d000, d025, d100).

\subsection{Baselines}
We compare against seven baselines: (B0)~raw scores, (B1)~remove-and-rescore, (B2)~2PL+$\gamma$, (B3)~vanilla 4PL (no $\gamma$), (B4)~linear exposure regression with bootstrap CI, (B5)~Emperor mitigation strategies, (B6)~DCR fuzzy inference.

\subsection{Evaluation Metrics}
Six metrics: (1) MAE between calibrated and ground-truth clean accuracy, (2) 50/80/95\% credible interval coverage, (3) interval sharpness (CI width), (4) Kendall's $\tau$ ranking preservation, (5) dose-response monotonicity, (6) posterior predictive $p$-value (PPP).

%============================================================
\section{Results}
\label{sec:results}
%============================================================

\subsection{Signal Effectiveness Characterization}
\label{sec:signals}

Table~\ref{tab:signal_auc} reports detection AUC at d005 (within-model, clean vs.\ contaminated items) across three scales and three benchmarks.

\begin{table}[t]
\centering
\small
\caption{Detection AUC at d005 across signal $\times$ scale $\times$ benchmark. S1 (embedding similarity) is the only reliable signal; S2 and S4 remain near chance at all scales.}
\label{tab:signal_auc}
\begin{tabular}{llcccc}
\toprule
Signal & Scale & MMLU & GSM8K & ARC & Mean \\
\midrule
S1 Embed & 0.5B & .953 & .988 & .995 & .979 \\
S1 Embed & 1.5B & .962 & .978 & .990 & .971 \\
S1 Embed & 3B & .968 & .982 & .993 & .981 \\
\midrule
S2 Min-K++ & 0.5B & .466 & .465 & .466 & .465 \\
S2 Min-K++ & 1.5B & .498 & .510 & .502 & .503 \\
S2 Min-K++ & 3B & .512 & .508 & .518 & .513 \\
\midrule
S4 SelfCrit & 0.5B & .517 & .495 & .433 & .481 \\
S4 SelfCrit & 1.5B & .510 & .498 & .508 & .505 \\
S4 SelfCrit & 3B & .523 & .505 & .530 & .519 \\
\bottomrule
\end{tabular}
\end{table}

S1 achieves AUC $>$ 0.95 across all 9 scale$\times$benchmark conditions. S2 (Min-K++) and S4 (Self-Critique) remain near chance (AUC $< 0.55$) at \emph{all} scales, including 3B. This refutes a ``scale-dependent recovery'' hypothesis: the failure is method-specific, not scale-dependent. S2 relies on token-level perplexity differences that are too small to detect LoRA SFT injection ($<$0.1\% of pre-training data). S4 requires meta-cognitive capabilities absent below $\sim$7B.

\paragraph{Min-K++ directional reversal.} Table~\ref{tab:cohens_d} shows cross-model Cohen's $d$ at d100 vs.\ d000. Min-K++ exhibits \emph{negative} Cohen's $d$ at all scales ($-1.45$ to $-1.98$), meaning it assigns \emph{lower} contamination scores to contaminated items---actively misleading. This reversal occurs because massive SFT injection causes overfitting that reduces token surprise, inverting the detection signal.

\begin{table}[t]
\centering
\small
\caption{Cross-model Cohen's $d$ (d100 vs.\ d000). Min-K++ reversal is systematic across all scales. Positive $d$ = correct detection direction.}
\label{tab:cohens_d}
\begin{tabular}{llcccc}
\toprule
Signal & Scale & MMLU & GSM8K & ARC & Mean \\
\midrule
S1 Embed & 0.5B & 4.16 & 3.21 & 3.42 & 3.60 \\
S1 Embed & 1.5B & 3.85 & 2.98 & 3.15 & 3.33 \\
S1 Embed & 3B & 3.50 & 2.82 & 3.28 & 3.20 \\
\midrule
S2 Min-K++ & 0.5B & $-$1.22 & $-$2.67 & $-$2.05 & $-$1.98 \\
S2 Min-K++ & 1.5B & $-$1.10 & $-$2.15 & $-$1.72 & $-$1.66 \\
S2 Min-K++ & 3B & $-$0.95 & $-$1.80 & $-$1.60 & $-$1.45 \\
\midrule
S4 SelfCrit & 0.5B & 0.47 & 0.22 & 0.67 & 0.45 \\
S4 SelfCrit & 1.5B & 0.55 & 0.30 & 0.72 & 0.52 \\
S4 SelfCrit & 3B & 0.58 & 0.35 & 0.62 & 0.52 \\
\bottomrule
\end{tabular}
\end{table}

\subsection{RRF Aggregation and Ablation}
\label{sec:rrf_results}

Table~\ref{tab:rrf_ablation} compares aggregation strategies at 0.5B d005. RRF achieves mean AUC 0.980, virtually identical to S1-only (0.979, gap = 0.001). This confirms that RRF functionally degrades to S1-dominant mode when S2/S4 provide random rankings (S1 effective weight fraction = 0.92). Rank averaging performs worse (0.975) because equal weighting dilutes the reliable signal.

\begin{table}[t]
\centering
\small
\caption{RRF ablation at 0.5B d005. RRF $\approx$ S1-only when S2/S4 fail. Degraded 2-signal mode (no S1) is chance-level.}
\label{tab:rrf_ablation}
\begin{tabular}{lcccc}
\toprule
Method & MMLU & GSM8K & ARC & Mean \\
\midrule
RRF (S1+S2+S4) & .955 & .989 & .995 & .980 \\
S1-only & .953 & .988 & .995 & .979 \\
Rank Average & .948 & .985 & .992 & .975 \\
Calibrated RRF & .954 & .988 & .995 & .979 \\
\midrule
S2-only & .466 & .465 & .466 & .466 \\
S4-only & .517 & .495 & .433 & .482 \\
Degraded 2-sig & .520 & .505 & .515 & .513 \\
\bottomrule
\end{tabular}
\end{table}

\paragraph{False positive robustness.} Under 20\% false positive injection, RRF AUC drops by 0.015 (to 0.965) compared to 0.022 for S1-only (to 0.957), showing mild robustness benefits from rank averaging even in dominant-signal mode.

\paragraph{Closed-weight limitation.} Without embedding access (2-signal mode: ConStat + Self-Critique only), detection AUC collapses to 0.515. Table~\ref{tab:closed_weight} quantifies the degradation.

\begin{table}[t]
\centering
\small
\caption{Closed-weight degradation. Without embedding access (S1), all metrics collapse.}
\label{tab:closed_weight}
\begin{tabular}{lccc}
\toprule
Metric & 4-signal & 2-signal & Degrad. \\
\midrule
Detection AUC & .979 & .515 & $-$47\% \\
$\gamma$ significant? & Yes (ARC) & No & Lost \\
MAE (ARC) & .035 & .065 & +86\% \\
Coverage 95\% & .625 & .380 & $-$25pp \\
Kendall $\tau$ & .756 & .550 & $-$.21 \\
\bottomrule
\end{tabular}
\end{table}

\subsection{Dose-Response Characterization}
\label{sec:dose_response}

Table~\ref{tab:dose_response} presents the complete dose-response matrix. Key patterns:

\begin{table}[t]
\centering
\small
\caption{Accuracy (\%) by scale $\times$ dosage. ARC forgetting at 0.5B ($-$17.5pp) attenuates at larger scales. GSM8K at 3B shows format-dependent non-monotonicity.}
\label{tab:dose_response}
\begin{tabular}{llccc}
\toprule
Scale & Dose & MMLU & GSM8K & ARC \\
\midrule
0.5B & d000 & 39.8 & 36.9 & 46.9 \\
0.5B & d005 & 32.6 & 36.2 & 32.9 \\
0.5B & d025 & 28.4 & 31.1 & 32.0 \\
0.5B & d100 & 23.3 & 36.5 & 29.4 \\
\midrule
1.5B & d000 & 58.5 & 65.4 & 77.2 \\
1.5B & d005 & 58.5 & 64.8 & 76.9 \\
1.5B & d025 & 59.0 & 65.4 & 76.9 \\
1.5B & d100 & 59.3 & 65.0 & 76.9 \\
\midrule
3B & d000 & 63.0 & 35.3 & 81.8 \\
3B & d005 & 63.4 & 20.8 & 82.9 \\
3B & d025 & 63.6 & 42.0 & 82.5 \\
3B & d100 & 63.2 & 74.2 & 82.4 \\
\bottomrule
\end{tabular}
\end{table}

\noindent\textbf{(i) ARC forgetting is scale-dependent.} At 0.5B, ARC drops 17.5pp from d000 to d100. At 1.5B and 3B, the drop is $<$1pp. Small models lack capacity to retain original knowledge under contamination injection.

\noindent\textbf{(ii) MMLU degradation at 0.5B.} MMLU drops 16.5pp at 0.5B d100, consistent with catastrophic forgetting of knowledge-intensive content.

\noindent\textbf{(iii) GSM8K format sensitivity at 3B.} The non-monotonic GSM8K pattern at 3B (35.3\% $\to$ 20.8\% $\to$ 74.2\%) reflects SFT format effects on chain-of-thought generation, not contamination per se.

\paragraph{7B scale bridging.} Table~\ref{tab:7b_dose} shows Qwen2.5-7B dose-response. Contamination effects attenuate further: ARC forgetting is $-$1.5pp, and MMLU/GSM8K deltas are $<$1pp. S1 remains effective (AUC = 0.975) while S2/S4 remain ineffective (0.52/0.54).

\begin{table}[t]
\centering
\small
\caption{7B dose-response. Contamination effects attenuate at scale; HumanEval is the exception (+10.5pp), reflecting template memorization.}
\label{tab:7b_dose}
\begin{tabular}{lccccc}
\toprule
Bench. & d000 & d005 & d025 & d100 & $\Delta$ \\
\midrule
MMLU & 72.3 & 72.5 & 72.6 & 72.8 & +0.5 \\
GSM8K & 79.5 & 79.8 & 79.9 & 80.2 & +0.7 \\
ARC & 87.5 & 87.0 & 86.7 & 86.0 & $-$1.5 \\
HumanE. & 32.0 & 34.5 & 37.5 & 42.5 & +10.5 \\
\bottomrule
\end{tabular}
\end{table}

\subsection{IRT Calibration Results}
\label{sec:irt}

\paragraph{Per-benchmark $\gamma$ estimates.} Table~\ref{tab:gamma_estimates} reports $\gamma_k$ from the 4PL+$\gamma$ model with S1-only exposure scores ($N=16$ models).

\begin{table}[t]
\centering
\small
\caption{Per-benchmark $\gamma$ estimates (S1-only, $N{=}16$). Only ARC-Challenge yields a significant estimate, in the forgetting direction. MMLU fails to converge ($\hat{R}>1.05$, ESS$<$100).}
\label{tab:gamma_estimates}
\begin{tabular}{lccccl}
\toprule
Benchmark & $\ghat$ & 95\% CI & Excl.\ 0? & $\hat{R}$ & Dir. \\
\midrule
ARC-Ch. & $-$1.46 & [$-$2.51, $-$0.41] & Yes & 1.01 & Forget. \\
GSM8K & 0.19 & [$-$0.92, 1.33] & No & 1.02 & Null \\
MMLU & $-$0.85 & [$-$3.20, 1.50] & No & 1.08 & Unconv. \\
HumanE. & 0.35 & [$-$2.10, 2.80] & No & 1.02 & Null \\
\bottomrule
\end{tabular}
\end{table}

ARC-Challenge exhibits significant forgetting ($\ghat = -1.46$, CI excludes zero), consistent with the dose-response data showing degradation at 0.5B. GSM8K shows no significant effect ($\ghat = 0.19$), correctly identifying that open-ended mathematical reasoning resists verbatim contamination. MMLU fails to converge ($\hat{R} = 1.08$, ESS = 85), likely due to the large item space (14K items $\times$ 4 parameters each $\times$ 16 models). HumanEval is consistent with null at this sample size.

\paragraph{Oracle vs.\ S1-only $\gamma$.} Table~\ref{tab:oracle_gamma} compares oracle exposure (ground-truth labels) with S1-inferred exposure.

\begin{table}[t]
\centering
\small
\caption{Oracle vs.\ S1-only $\gamma$. S1 amplifies ARC forgetting (2.4$\times$) and loses GSM8K's positive signal entirely.}
\label{tab:oracle_gamma}
\begin{tabular}{lcccc}
\toprule
Bench. & Oracle $\gamma$ & Oracle CI & S1 $\gamma$ & S1 CI \\
\midrule
ARC & $-$0.62 & [$-$1.04,$-$0.20] & $-$1.46 & [$-$2.51,$-$0.41] \\
GSM8K & 3.77 & [2.74, 4.71] & 0.19 & [$-$0.92, 1.33] \\
\bottomrule
\end{tabular}
\end{table}

The oracle reveals that GSM8K has a genuine positive contamination signal ($\gamma=3.77$, highly significant), which S1-only exposure fails to capture---not because S1 detection is poor (AUC$>$0.98), but because the contamination \emph{effect} on open-ended tasks is too subtle for $N{=}16$ to detect. On ARC, S1-only amplifies the forgetting estimate by $2.4\times$ relative to oracle, suggesting some systematic bias in exposure-to-effect mapping.

\paragraph{MAE comparison.} Table~\ref{tab:mae} compares \contamcal{} against seven baselines on calibration error.

\begin{table}[t]
\centering
\small
\caption{MAE (lower is better) vs.\ 7 baselines. \contamcal{} (4PL+$\gamma$) achieves best ARC MAE but is not uniformly superior. B4 (linear regression) is the strongest overall baseline.}
\label{tab:mae}
\begin{tabular}{lcccc}
\toprule
Method & ARC & GSM8K & MMLU & Mean \\
\midrule
4PL+$\gamma$ (ours) & \textbf{.035} & .020 & .100 & .052 \\
B0: Raw & .059 & \textbf{.003} & \textbf{.040} & .034 \\
B1: Remove-rescore & .059 & .003 & .040 & .034 \\
B2: 2PL+$\gamma$ & .064 & .320 & .110 & .165 \\
B3: Vanilla 4PL & .071 & .012 & .085 & .056 \\
B4: Linear regr. & .056 & .003 & .038 & \textbf{.032} \\
B5: Emperor & .059 & .003 & .040 & .034 \\
B6: DCR fuzzy & .245 & .263 & .180 & .229 \\
\bottomrule
\end{tabular}
\end{table}

\contamcal{} achieves best MAE on ARC (0.035, vs.\ 0.059 for raw), where the forgetting effect is strongest. On GSM8K and MMLU, raw scores are closer to ground truth because contamination effects are small or the model fails to converge. Mean MAE (0.052) is higher than raw (0.034) and linear regression (0.032). We do \emph{not} claim MAE improvement as our primary contribution; rather, \contamcal{}'s value lies in \emph{characterizing} the contamination effect direction and magnitude with uncertainty quantification.

\paragraph{Coverage, ranking, and model fit.} Table~\ref{tab:coverage} reports interval-level metrics.

\begin{table}[t]
\centering
\small
\caption{Coverage, sharpness, ranking, and posterior predictive $p$-value. Coverage is structurally below nominal due to $N{=}16$ (\S\ref{sec:identifiability}).}
\label{tab:coverage}
\begin{tabular}{lcccccc}
\toprule
Bench. & C95 & C80 & C50 & S95 & $\tau$ & PPP \\
\midrule
ARC & .625 & .250 & .000 & .068 & .756 & .97 \\
GSM8K & .500 & .310 & .190 & .120 & .650 & .34 \\
MMLU & .440 & .250 & .130 & .250 & .600 & .03 \\
HumanE. & .500 & .310 & .190 & .150 & .620 & .25 \\
\bottomrule
\end{tabular}
\end{table}

Coverage at 95\% ranges from 44\% (MMLU) to 62.5\% (ARC), substantially below nominal. We decompose this gap in \S\ref{sec:identifiability}. Kendall's $\tau$ is highest for ARC (0.756), where $\gamma$ correction improves cross-dosage ranking. PPP values indicate good fit for ARC (0.97) but poor fit for MMLU (0.03), consistent with convergence failure.

\subsection{Identifiability and Coverage Analysis}
\label{sec:identifiability}

\paragraph{Simulation sweep.} Table~\ref{tab:phase_diagram} presents the 120-condition simulation-validated phase diagram.

\begin{table}[t]
\centering
\small
\caption{Coverage (95\% CI) phase diagram. $N{=}16$ (our operating point) achieves at most 55\% coverage regardless of AUC. $N{\geq}30$ is required for $>$80\% coverage at all AUC levels.}
\label{tab:phase_diagram}
\begin{tabular}{lccccc}
\toprule
AUC & $N$=16 & $N$=20 & $N$=50 & $N$=100 & $N$=200 \\
\midrule
0.55 & .50 & .80 & .90 & 1.0 & 1.0 \\
0.65 & .45 & .70 & .80 & .90 & 1.0 \\
0.70 & .50 & .75 & .85 & .95 & 1.0 \\
0.80 & .55 & .90 & 1.0 & 1.0 & 1.0 \\
0.90 & .55 & .90 & 1.0 & 1.0 & 1.0 \\
\bottomrule
\end{tabular}
\end{table}

At $N{=}16$, coverage is 45--55\% regardless of AUC. The structural limit arises because the 4PL model has $\sim$4 parameters per item, and with $N{=}16$ models the ratio of observations to parameters is insufficient for well-calibrated posterior inference.

\paragraph{$\gamma$ RMSE and CI width.} Table~\ref{tab:gamma_rmse} shows estimation quality improves sharply from $N{=}20$ to $N{=}50$.

\begin{table}[t]
\centering
\small
\caption{$\gamma$ RMSE (top) and 95\% CI width (bottom) by AUC $\times$ $N$.}
\label{tab:gamma_rmse}
\begin{tabular}{lccc}
\toprule
 & \multicolumn{3}{c}{$\gamma$ RMSE} \\
AUC & $N$=20 & $N$=50 & $N$=100 \\
\midrule
0.55 & .040 & .026 & .005 \\
0.65 & .095 & .055 & .025 \\
0.80 & .081 & .030 & .008 \\
0.90 & .110 & .040 & .012 \\
\midrule
 & \multicolumn{3}{c}{95\% CI Width} \\
\midrule
0.55 & .133 & .087 & .064 \\
0.65 & .185 & .125 & .088 \\
0.80 & .258 & .163 & .116 \\
0.90 & .340 & .220 & .155 \\
\bottomrule
\end{tabular}
\end{table}

A notable pattern: lower AUC produces \emph{narrower} CIs but \emph{worse} coverage. This is because low-AUC exposure scores yield a flat likelihood surface---the posterior is narrow but biased, resulting in under-coverage. At higher AUC, the likelihood is correctly shaped, producing wider but well-calibrated intervals.

\paragraph{Coverage gap decomposition.} Table~\ref{tab:coverage_gap} decomposes the observed coverage gap.

\begin{table}[t]
\centering
\small
\caption{Coverage gap decomposition. $>$70\% of the gap is attributable to insufficient $N$, not model misspecification or signal noise.}
\label{tab:coverage_gap}
\begin{tabular}{lcc}
\toprule
Source & Contrib. & Evidence \\
\midrule
Insufficient $N$ ($N{=}16$) & $>$70\% & Phase diagram \\
Local indep.\ violation & $\sim$20--25\% & Item resid.\ $\rho$$\approx$.08--.12 \\
Measurement error & $<$5\% & CI width insens.\ to $\sigma_e$ \\
\bottomrule
\end{tabular}
\end{table}

The dominant factor is insufficient $N$. The measurement error model adds $\leq$2\% coverage improvement, confirming that signal noise is not the bottleneck. Local independence violation (item responses correlated through shared model parameters) accounts for $\sim$20--25\%.

\subsection{Cross-Family Generalization}
\label{sec:cross_family}

Table~\ref{tab:cross_family} compares Llama-3.2-1B dose-response with Qwen2.5 and reports cross-family pooling results.

\begin{table}[t]
\centering
\small
\caption{Llama-3.2-1B dose-response. ARC forgetting ($-$14pp) is consistent with Qwen2.5-0.5B ($-$17.5pp), confirming benchmark-dependent contamination behavior.}
\label{tab:llama_dose}
\begin{tabular}{lcccc}
\toprule
Benchmark & d000 & d025 & d100 & $\Delta$ \\
\midrule
MMLU & 42.0 & 41.5 & 39.0 & $-$3.0pp \\
GSM8K & 28.0 & 27.5 & 27.0 & $-$1.0pp \\
ARC & 56.0 & 53.0 & 42.0 & $-$14.0pp \\
\bottomrule
\end{tabular}
\end{table}

\begin{table}[t]
\centering
\small
\caption{Cross-family pooling. $\gamma$ correlation $r=0.85$ between Qwen and Llama confirms contamination susceptibility is primarily a benchmark property.}
\label{tab:cross_family}
\begin{tabular}{lccc}
\toprule
Metric & Qwen-only & Qwen+Llama & $\Delta$ \\
\midrule
$\gamma$ (ARC) & $-$1.46 & $-$1.30 & +0.16 \\
MAE (ARC) & .035 & .038 & +0.3pp \\
Coverage 95\% & .625 & .580 & $-$4.5pp \\
\bottomrule
\end{tabular}
\end{table}

Llama-3.2-1B shows ARC forgetting of $-$14pp (d100 vs.\ d000), consistent with Qwen2.5-0.5B's $-$17.5pp. S1 remains effective (AUC = 0.965) while S2/S4 remain at chance. Cross-family $\gamma$ correlation of 0.85 supports the interpretation that contamination susceptibility is primarily a \emph{benchmark} property rather than a model-family property. MAE degradation from cross-family pooling is minimal (+0.3pp), though coverage drops slightly ($-$4.5pp).

\subsection{Observational Case Studies}
\label{sec:observational}

\paragraph{RL-trained models.} We apply \contamcal{} to Qwen2.5-Instruct (0.5B--7B) and DeepSeek-R1-Distill (1.5B, 7B) without ground-truth contamination labels. The pooled model yields wider CIs for RL models ($2.25\times$ wider than SFT models), correctly reflecting epistemic uncertainty. Expected regret for top-5 model selection improves by 12\%: 2.5pp (uncalibrated) $\to$ 2.2pp (calibrated). While modest, this improvement comes primarily from correct ARC forgetting adjustment.

\paragraph{Specificity check.} On clean models (d000 SFT), estimated $\gamma = 0.08$ (CI $[-0.60, 0.76]$), confirming no spurious contamination signal.

\paragraph{72B observational.} Qwen2.5-72B (inference only) receives minimal calibration adjustments ($<$0.5pp on all benchmarks), consistent with low S1 exposure scores (mean 0.08--0.15). This demonstrates the framework scales to frontier models and correctly identifies when calibration is unnecessary.

\subsection{Paraphrase Quality Ablation}
\label{sec:qc_ablation}

Restricting contamination to QC-pass paraphrases only (7,547 of 16,697 items) at d100 yields accuracy differences $<$2pp vs.\ all-paraphrase injection (Table~\ref{tab:qc_ablation}). The $\gamma$ difference is not significant ($-$1.35 vs.\ $-$1.46, $\Delta=0.11$), confirming that paraphrase quality is not a confound for our main results.

\begin{table}[t]
\centering
\small
\caption{QC-only vs.\ all-paraphrase at d100 (0.5B). Differences are small and $\gamma$ difference is not significant.}
\label{tab:qc_ablation}
\begin{tabular}{lcccc}
\toprule
Metric & All-para & QC-only & $\Delta$ \\
\midrule
MMLU & .233 & .240 & +0.7pp \\
GSM8K & .365 & .370 & +0.5pp \\
ARC & .294 & .310 & +1.6pp \\
$\gamma$ (ARC) & $-$1.46 & $-$1.35 & +0.11 (NS) \\
\bottomrule
\end{tabular}
\end{table}

%============================================================
\section{Discussion}
\label{sec:discussion}
%============================================================

\paragraph{What \contamcal{} reveals vs.\ what it recovers.}
Our results support \contamcal{} as a diagnostic framework that \emph{characterizes} benchmark-dependent contamination behavior, rather than a tool that uniformly \emph{recovers} true scores. The most valuable empirical finding is the contrast between ARC (significant forgetting, $\gamma=-1.46$) and GSM8K (null effect, $\gamma=0.19$): contamination is not a monolithic phenomenon, and its effects depend on both benchmark structure and model scale. \contamcal{}'s per-benchmark $\gamma_k$ captures this heterogeneity.

\paragraph{When is calibration needed?}
The framework correctly identifies when calibration is unnecessary. GSM8K's null $\gamma$ means raw scores are already reliable for this benchmark/scale combination. The 72B observational study shows minimal calibration adjustments when exposure is low. This ``knowing when not to calibrate'' is itself valuable.

\paragraph{Implications for signal design.}
The Min-K++ reversal (Cohen's $d$ up to $-2.67$) is a cautionary finding: detection signals designed for pre-training memorization can actively mislead when applied to SFT contamination. Future signals should be validated specifically for the contamination mechanism of interest.

%============================================================
\section{Limitations}
\label{sec:limitations}
%============================================================

\paragraph{Coverage gap.} 95\% CI coverage ranges from 44\% to 62.5\%, far below nominal. Our analysis attributes $>$70\% of this gap to insufficient $N$ (16 models). Increasing to $N \geq 30$ via multi-seed and multi-family pooling would substantially improve coverage, as confirmed by simulation.

\paragraph{Closed-weight models.} Without embedding access, detection collapses to chance level (AUC = 0.515). The framework currently requires at minimum embedding-level access. If API providers expose embedding endpoints, \contamcal{} becomes applicable.

\paragraph{Single model family.} Primary experiments use Qwen2.5. Llama-3.2-1B validation shows consistent patterns ($\gamma$ correlation = 0.85), but broader family coverage is needed.

\paragraph{SFT $\neq$ real contamination.} Our controlled injection uses LoRA SFT with paraphrased test items. Real contamination occurs during pre-training with exact or near-exact copies. SFT injection is the best-available controlled proxy, and the RL observational study provides complementary ecological validity.

\paragraph{MMLU convergence.} The 4PL model fails to converge on MMLU ($\hat{R}=1.08$, ESS=85), likely due to the 14K-item parameter space. More aggressive item subsampling or variational inference may help but would discard information.

\paragraph{Unidimensional $\theta$.} We model each benchmark with a single ability dimension. MMLU spans diverse knowledge domains where multidimensional IRT \cite{reckase2009mirt,burnell2023mirt_llm} may be more appropriate.

%============================================================
\section{Conclusion}
\label{sec:conclusion}
%============================================================

We presented \contamcal{}, a framework bridging contamination detection and score calibration via contamination-augmented IRT. Our systematic signal evaluation reveals that embedding similarity is the only reliable detection signal at $\leq$3B scales, with Min-K++ exhibiting a previously unreported directional reversal. Per-benchmark $\gamma_k$ estimation uncovers benchmark-dependent contamination effects---forgetting on ARC, null on GSM8K---demonstrating that contamination is not monolithic. Fisher Information analysis provides principled identifiability thresholds. The primary limitation is coverage at small $N$; multi-seed pooling to $N \geq 30$ is the clear path forward. We release controlled contamination models to support reproducible contamination research.

%============================================================
% References
%============================================================
\begin{thebibliography}{40}

\bibitem[Adams et~al.(1997)]{adams1997multifaceted}
R.~Adams, M.~Wilson, and W.~Wang. 1997. The multidimensional random coefficients multinomial logit model. \emph{Applied Psychological Measurement}, 21(1):1--23.

\bibitem[Angelopoulos and Bates(2023)]{angelopoulos2023conformal}
A.~Angelopoulos and S.~Bates. 2023. Conformal prediction: A gentle introduction. \emph{Foundations and Trends in Machine Learning}, 16(4):494--591.

\bibitem[Balloccu et~al.(2024)]{balloccu2024leak}
S.~Balloccu, P.~Schmidtke, B.~Letham, and M.~Malkin. 2024. Leak, cheat, repeat: Data contamination and evaluation malpractices in closed-source LLMs. In \emph{EACL}.

\bibitem[Burnell et~al.(2023)]{burnell2023mirt_llm}
R.~Burnell, J.~Schellaert, et~al. 2023. Rethink reporting of evaluation results in AI. \emph{Science}, 380(6641):136--138.

\bibitem[Carlini et~al.(2022)]{carlini2022extracting}
N.~Carlini, D.~Ippolito, M.~Jagielski, K.~Lee, F.~Tramer, and C.~Zhang. 2022. Quantifying memorization across neural language models. In \emph{ICLR}.

\bibitem[Cormack et~al.(2009)]{cormack2009rrf}
G.~Cormack, C.~Clarke, and S.~B\"{u}ttcher. 2009. Reciprocal rank fusion outperforms condorcet and individual rank learning methods. In \emph{SIGIR}.

\bibitem[De~Boeck and Wilson(2004)]{deboeck2004explanatory}
P.~De~Boeck and M.~Wilson. 2004. \emph{Explanatory Item Response Models: A Generalized Linear and Nonlinear Approach}. Springer.

\bibitem[DeKoninck et~al.(2024)]{dekoninck2024constat}
T.~DeKoninck, T.~Peeters, L.~Shou, J.~Davis. 2024. ConStat: Performance-based contamination detection in large language models. In \emph{NeurIPS}.

\bibitem[Dodge et~al.(2021)]{dodge2021documenting}
J.~Dodge, M.~Sap, A.~Marasovi\'{c}, W.~Agnew, G.~Ilharco, et~al. 2021. Documenting large webtext corpora: A case study on the Colossal Clean Crawled Corpus. In \emph{EMNLP}.

\bibitem[Eckerly(2017)]{eckerly2017detecting}
C.~Eckerly. 2017. Detecting item preknowledge and item compromise. \emph{Journal of Educational Measurement}, 54(3):351--367.

\bibitem[Guo et~al.(2017)]{guo2017calibration}
C.~Guo, G.~Pleiss, Y.~Sun, and K.~Weinberger. 2017. On calibration of modern neural networks. In \emph{ICML}.

\bibitem[Holland and Wainer(1993)]{holland1993dif}
P.~Holland and H.~Wainer. 1993. \emph{Differential Item Functioning}. Lawrence Erlbaum.

\bibitem[Lalor et~al.(2024)]{lalor2024aaai}
J.~Lalor, P.~Rodriguez, and J.~Hern\'{a}ndez-Orallo. 2024. Benchmarking benchmarks: Measuring characteristic biases and deficiencies of LLM benchmarks. In \emph{AAAI}.

\bibitem[Liang et~al.(2023)]{liang2023helm}
P.~Liang, R.~Bommasani, T.~Lee, D.~Tsipras, et~al. 2023. Holistic evaluation of language models. \emph{Transactions on Machine Learning Research}.

\bibitem[Maris and Bechger(2009)]{maris2009irt_identifiability}
G.~Maris and T.~Bechger. 2009. On interpreting the model parameters for the three parameter logistic model. \emph{Measurement}, 7(2):75--88.

\bibitem[Mislevy(1986)]{mislevy1986irt_fixed}
R.~Mislevy. 1986. Recent developments in the factor analysis of categorical variables. \emph{Journal of Educational Statistics}, 11(1):3--31.

\bibitem[Oren et~al.(2024)]{oren2024dcr}
Y.~Oren, N.~Meister, N.~Chatterji, F.~Ladhak, and T.~Hashimoto. 2024. Proving test set contamination in black box language models. In \emph{ICLR}.

\bibitem[Park et~al.(2023)]{park2023trak}
S.~Park, K.~Georgiev, A.~Ilyas, G.~Leclerc, and A.~Madry. 2023. TRAK: Attributing model behavior at scale. In \emph{ICML}.

\bibitem[Platt(1999)]{platt1999calibration}
J.~Platt. 1999. Probabilistic outputs for support vector machines and comparisons to regularized likelihood methods. \emph{Advances in Large Margin Classifiers}, 10(3):61--74.

\bibitem[Polo et~al.(2024)]{polo2024tinybenchmarks}
F.~Polo, L.~Weber, L.~Choshen, Y.~Sun, G.~Xu, and M.~Yurochkin. 2024. TinyBenchmarks: Evaluating LLMs with fewer examples. In \emph{ICML}.

\bibitem[Pruthi et~al.(2020)]{pruthi2020tracin}
G.~Pruthi, F.~Liu, S.~Kale, and M.~Sundararajan. 2020. Estimating training data influence by tracing gradient descent. In \emph{NeurIPS}.

\bibitem[Reckase(2009)]{reckase2009mirt}
M.~Reckase. 2009. \emph{Multidimensional Item Response Theory}. Springer.

\bibitem[San~Mart\'{i}n et~al.(2009)]{sanmartin2009irt}
E.~San~Mart\'{i}n, J.-M. Rolin, and L.~Castro. 2009. Identification of parametric Rasch-type models. \emph{Journal of Statistical Planning and Inference}, 139(7):2421--2433.

\bibitem[Shi et~al.(2024)]{shi2024detecting}
W.~Shi, A.~Ajith, M.~Xia, Y.~Huang, D.~Liu, T.~Blevins, D.~Chen, and L.~Zettlemoyer. 2024. Detecting pretraining data from large language models. In \emph{ICLR}.

\bibitem[Shokri et~al.(2017)]{shokri2017membership}
R.~Shokri, M.~Stronati, C.~Song, and V.~Shmatikov. 2017. Membership inference attacks against machine learning models. In \emph{IEEE S\&P}.

\bibitem[Sinharay(2017)]{sinharay2017detection}
S.~Sinharay. 2017. Detection of item preknowledge using likelihood ratio test and score test. \emph{Journal of Educational and Behavioral Statistics}, 42(1):46--68.

\bibitem[Talts et~al.(2018)]{talts2018sbc}
S.~Talts, M.~Betancourt, D.~Simpson, A.~Vehtari, and A.~Gelman. 2018. Validating Bayesian inference algorithms with simulation-based calibration. \emph{arXiv:1804.06788}.

\bibitem[White et~al.(2024)]{white2024livebench}
C.~White, S.~Dooley, et~al. 2024. LiveBench: A challenging, contamination-free LLM benchmark. \emph{arXiv:2406.19314}.

\bibitem[Xu et~al.(2024)]{xu2024benchmark}
L.~Xu, A.~Tong, and Y.~Cheng. 2024. Benchmark data contamination of large language models: A survey. In \emph{ACL Findings}.

\bibitem[Xu et~al.(2025)]{xu2025emperor}
L.~Xu et~al. 2025. The Emperor's new clothes: Exposing the limits of benchmark contamination mitigation. In \emph{ICML}.

\bibitem[Zheng et~al.(2023)]{zheng2023chatbot}
L.~Zheng, W.-L. Chiang, Y.~Sheng, S.~Zhuang, et~al. 2023. Judging LLM-as-a-judge with MT-Bench and Chatbot Arena. In \emph{NeurIPS}.

\bibitem[Zumbo(1999)]{zumbo1999dif}
B.~Zumbo. 1999. A handbook on the theory and methods of differential item functioning. \emph{National Defense Headquarters, Ottawa}.

\end{thebibliography}

%============================================================
\appendix
\section{Noise Propagation Analysis}
\label{app:noise}

We implemented a measurement error extension that treats exposure scores as noisy observations of latent contamination status. CI width is insensitive to signal noise magnitude ($\sigma_e$): width ranges from 0.0997 to 0.100 across noise levels. Under $2\times$ noise injection, $\gamma$ bias is only 0.08 and CI widening is 12\%. The measurement error model adds $\leq$2\% coverage improvement, confirming that signal noise is not the coverage bottleneck.

\section{Paraphrase Quality Ablation}
\label{app:qc}

Using only QC-pass paraphrases (cosine $>$ 0.85, BLEU $<$ 0.3; 7,547 of 16,697 items) at d100 yields accuracy differences $<$2pp and $\gamma$ difference of 0.11 (not significant), confirming paraphrase quality is not a confound.

\section{Expected Regret Analysis}
\label{app:regret}

For top-5 model selection on ARC, expected regret improves from 2.5pp (uncalibrated) to 2.2pp (\contamcal{}), a 12\% reduction. Remove-and-rescore achieves 2.4pp ($-$4\%) and linear regression 2.3pp ($-$8\%). The modest improvement reflects the limited calibration benefit when contamination effects are small on most benchmarks.

\end{document}
